% =============================================================================
% demo-performance.tex — Demo document for the swarmperf theme
%
% Compile: pdflatex demo-performance.tex  (or lualatex / xelatex)
%          (run twice for TOC and cross-references)
% =============================================================================

\documentclass[11pt, a4paper]{article}

\usepackage{swarmperf}

% ── Metadata ─────────────────────────────────────────────────────────────────

\title{SwarmPerf Theme Demo}
\subtitle{A fast, minimal showcase of the Performance Theme}
\author{LaTeX Helper Swarm Project}
\date{May 14, 2026}

% ── Bibliography (optional — uncomment if you have a .bib file) ─────────────
% \usepackage[backend=biber, style=numeric]{biblatex}
% \addbibresource{references.bib}

\begin{document}

% ══════════════════════════════════════════════════════════════════════════════
% Title Page
% ══════════════════════════════════════════════════════════════════════════════

\swarmtitlepage

% ══════════════════════════════════════════════════════════════════════════════
% Table of Contents
% ══════════════════════════════════════════════════════════════════════════════

\tableofcontents
\newpage

% ══════════════════════════════════════════════════════════════════════════════
% Section 1: Introduction
% ══════════════════════════════════════════════════════════════════════════════

\section{Introduction}

This document demonstrates the \code{swarmperf} theme, a minimal, fast-to-compile
LaTeX document style designed for speed and simplicity. Unlike the full-featured
\code{swarmbeauty} theme, \code{swarmperf} deliberately excludes heavy packages
such as TikZ, fontspec, microtype, and tcolorbox to minimize compilation time.

The design philosophy is simple: every package must earn its place. The result is
a theme that compiles \hltext{2--5x faster} than typical modern LaTeX setups while
still producing clean, professional documents suitable for technical reports,
academic papers, and internal documentation.

\subsection{Motivation}

Many LaTeX documents do not need elaborate decorations, syntax-highlighted code
blocks via external tools, or custom OpenType fonts. For these use cases, the
overhead of loading fontspec, minted, tcolorbox, and TikZ can add significant
compilation time---especially on CI/CD servers or in iterative edit-compile cycles.

The \code{swarmperf} theme provides sensible defaults for clean document
formatting without the weight.

\begin{noteblock}
This theme works with \textbf{pdfLaTeX}, \textbf{XeLaTeX}, and
\textbf{LuaLaTeX}. No \code{-{}-shell-escape} flag is required.
\end{noteblock}

% ══════════════════════════════════════════════════════════════════════════════
% Section 2: Typography & Text
% ══════════════════════════════════════════════════════════════════════════════

\section{Typography and Text}

The theme uses the default Computer Modern fonts (or Latin Modern if available).
Text is set with sans-serif headings and serif body text. Mathematics uses the
standard math fonts.

\subsection{Inline Elements}

You can highlight \hltext{important terms} using the \code{\\hltext} command.
For inline code, use \code{\\code}: \codeesc{some_function()}.
For code with special characters: \codeesc{text_with_underscores}.

\subsection{Lists}

\begin{itemize}
  \item First item in an unordered list
  \item Second item with \hltext{emphasis}
  \item Third item with nested content:
    \begin{itemize}
      \item Nested item A
      \item Nested item B
    \end{itemize}
\end{itemize}

\begin{enumerate}
  \item Step one: load the performance theme
  \item Step two: write your content
  \item Step three: compile with any LaTeX engine
\end{enumerate}

% ══════════════════════════════════════════════════════════════════════════════
% Section 3: Block Environments
% ══════════════════════════════════════════════════════════════════════════════

\section{Block Environments}

The theme provides lightweight block environments for notes, tips, and warnings.
These use basic minipage and color commands---no tcolorbox overhead.

\begin{noteblock}
This is a \textbf{note} block. Use it for general observations, additional
context, or supplementary information that is relevant but not critical to
the main argument.
\end{noteblock}

\begin{tipblock}
This is a \textbf{tip} block. Use it to highlight best practices, shortcuts,
or recommendations that can improve the reader's workflow or understanding.
\end{tipblock}

\begin{warningblock}
This is a \textbf{warning} block. Use it to alert the reader about potential
pitfalls, common mistakes, or deprecated features that they should be aware of.
\end{warningblock}

% ══════════════════════════════════════════════════════════════════════════════
% Section 4: Tables
% ══════════════════════════════════════════════════════════════════════════════

\section{Tables}

The theme integrates with \code{booktabs} for clean, professional tables.

\subsection{Booktabs Table}

\begin{table}[H]
  \centering
  \caption{Comparison of LaTeX compilation engines}
  \label{tab:engines}
  \begin{tabular}{l c c c}
    \swarmtoprule
    \textbf{Engine}   & \textbf{Speed} & \textbf{Font Support} & \textbf{Unicode} \\
    \swarmmidrule
    pdfLaTeX           & Fast           & Limited (Type 1)      & Partial          \\
    XeLaTeX            & Medium         & Full (OpenType)        & Full             \\
    LuaLaTeX           & Slow           & Full (OpenType)        & Full             \\
    \swarmbottomrule
  \end{tabular}
\end{table}

% ══════════════════════════════════════════════════════════════════════════════
% Section 5: Code Listings
% ══════════════════════════════════════════════════════════════════════════════

\section{Code Listings}

The \code{swarmperf} theme uses the \code{listings} package for code highlighting.
Unlike \code{minted}, \code{listings} requires no external dependencies (no
Pygments, no \code{-{}-shell-escape}).

\begin{codelisting}[Python]{Fibonacci sequence generator}
def fibonacci(n: int) -> list[int]:
    """Generate the first n Fibonacci numbers."""
    seq = [0, 1]
    for i in range(2, n):
        seq.append(seq[-1] + seq[-2])
    return seq[:n]

if __name__ == "__main__":
    result = fibonacci(10)
    print(f"First 10: {result}")
\end{codelisting}

\begin{codelisting}[TeX]{A simple equation}
\begin{equation}
  E = mc^2
  \label{eq:einstein}
\end{equation}
\end{codelisting}

% ══════════════════════════════════════════════════════════════════════════════
% Section 6: Mathematics
% ══════════════════════════════════════════════════════════════════════════════

\section{Mathematics}

The theme provides standard theorem environments with automatic numbering.

\begin{theorem}{Euler's Identity}{thm:euler}
  For any real number $x$, the following identity holds:
  \[
    e^{i\pi} + 1 = 0
  \]
\end{theorem}

\begin{definition}{Metric Space}{def:metric}
  A \emph{metric space} is a pair $(X, d)$ where $X$ is a set and
  $d\colon X \times X \to \mathbb{R}_{\geq 0}$ satisfies:
  \begin{enumerate}
    \item $d(x, y) = 0 \iff x = y$ (identity of indiscernibles)
    \item $d(x, y) = d(y, x)$ (symmetry)
    \item $d(x, z) \leq d(x, y) + d(y, z)$ (triangle inequality)
  \end{enumerate}
\end{definition}

\begin{lemma}{Nested Intervals}{lem:nested}
  Let $\{[a_n, b_n]\}_{n=1}^{\infty}$ be a sequence of closed, bounded
  intervals such that $[a_{n+1}, b_{n+1}] \subseteq [a_n, b_n]$ for all $n$.
  Then $\bigcap_{n=1}^{\infty} [a_n, b_n] \neq \emptyset$.
\end{lemma}

\subsection{Inline and Display Math}

The quadratic formula: $x = \frac{-b \pm \sqrt{b^2 - 4ac}}{2a}$

A more complex display equation:
\[
  \int_{-\infty}^{\infty} e^{-x^2} \, dx = \sqrt{\pi}
\]

% ══════════════════════════════════════════════════════════════════════════════
% Section 7: Summary
% ══════════════════════════════════════════════════════════════════════════════

\section{Summary}

This demo has showcased the following features of the \code{swarmperf} theme:

\begin{enumerate}
  \item \textbf{Title page} with colored accent rule
  \item \textbf{Section headings} with clean sans-serif formatting
  \item \textbf{Block environments} (note, tip, warning)
  \item \textbf{Styled tables} with booktabs
  \item \textbf{Syntax-highlighted code} via listings (no external deps)
  \item \textbf{Theorem environments} with section numbering
  \item \textbf{Inline code and emphasis} commands
  \item \textbf{Headers and footers} with page numbers
\end{enumerate}

\colorrule

\begin{center}
  \small\textcolor{spGray}{%
    Generated with \texttt{swarmperf} theme --- LaTeX Helper Swarm Project%
  }
\end{center}

\end{document}
